% EGA V, English translation by P. Blass, 1987 (cote 24): archive cote 24, site page 58 (PDF page 59).
% typescript. Draft compared: none.
% AI-transcribed by Claude Opus 5.5 (claude-opus-5-5, adaptive thinking, effort high), Message Batches API, batch ega24b (msgbatch_01RGJJeHdR96U16ZWFuUjtDx), 2026-11-13; not yet checked by eye.
\scanpage{24/58}
\mspage{58}

We shall now examine more precisely the case of surfaces (the
case of curves does not arise evidently, from the point of view of
irreducibility of hyperplane sections).
NB  I noticed with fright that the quadric is not entitled to the
use [Fr il l'emploi que j'ai fait de mot
``reglé,'' à être appelée réglée!]*
This is in disagreement with our fathers and it would be necessary
to invent a more adequate word for the notion used here.

\uline{Proposition 8.18}

Let us suppose that k is algebraically closed, X is integral
(respectively integral and normal) of dimension $\geq 2$ and proper
over k, let T be a closed finite subset of X such that X - T is
smooth and let f/X - T be unramified.  In order that the set of $\xi$
$\in P^{\vee}$ such that $Y_{\xi}$ should be geometrically irreducible (respectively
geometrically integral) of dimension 1 should have a complement of
codimension $\geq 2$ it is necessary and sufficient that the following
two conditions should be satisfied:  a) For every $x \in T$ there
exists a hyperplane section $Y_{\xi}$ ($\xi \in P(k)$ ) passing through x of
dimension d - 1 and which is irreducible, b) X' = X - T is ``ruled''
(\uline{sic}) for f \uline{or} there exists a hyperplane section $Y'_{\xi}$ ($\xi \in P^{\vee}(k)$ ) of
X' which is of dimension d - 1 (\uline{non}??? singular and irreducible.
[illegible ask AG]

Let us first assume that X is geometrically normal.  We have
already seen then by 8.14  a) that we can find a closed \uline{subset} $Z^{1}$
of P of codim $\geq 2$ such that $\xi \in P-Z^{1}$ implies that $Y_{\xi}$ is separable
over $k(\xi)$ and of dimension d-1 for such a $\xi$, it amounts to the same

\margin{* ASK AG}
